\chapter{定价、校准与对冲：从数值方法到误差分布}
\label{ch:lower-derivatives-numerics}

\begin{itemize}
  % Generated from curriculum/manifest.json; do not edit by hand.
\item 直接先修：下册《随机分析与无套利：从二次变差到风险中性测度》。

\end{itemize}

复制得到价格方程之后，还要把它算出来。以下用一层树说明风险中性概率的来源，再把空间导数变成差分系数，最后讨论报价反演的敏感度。


\MFQTransition{所求解的定价问题}
上一章已由自融资复制推导 Black--Scholes PDE。本节沿用同一方程与终端支付；数值方法改变的是求解方式，而非另行假定价格期望。含连续股息率 $q$ 时，股票方向的一阶系数为 $(r-q)S$。树法、差分与模拟必须针对同一合约比较。

\section{同一价格，不同误差来源}

欧式看涨闭式价格为
\[
 C=S_0\Phi(d_1)-Ke^{-rT}\Phi(d_2),\qquad
 d_{1,2}=\frac{\log(S_0/K)+(r\pm\sigma^2/2)T}{\sigma\sqrt T}.
\]
在 $S_0=K=100,r=2\%,\sigma=20\%,T=1$ 下，解析值是 $8.916037$。它是当前模型中的参考值，不是真实市场的“正确价格”。

% Generated by tools/build_figures.py; edit figures/figure-specs.json instead.
\MFQFigure{tex/figures/generated/lower-dv-methods.pdf}{不同定价方法使用相同的合约与模型输入，却引入不同的近似误差；加密网格、增加样本及比较方法可以帮助理解这些误差。}{lower-dv-methods}


这些方法近似同一个模型价格，但误差来源不同：
\begin{center}
\begin{tabular}{lll}
\toprule
方法 & 首要误差 & 误差的考察方式\\
\midrule
闭式 & 参数/模型错设 & 边界、单位与极限情形\\
CRR 树 & 时间离散 & $p\in[0,1]$ 与步数敏感性\\
隐式 PDE & 空间、时间、截断 & 边界条件与三维网格敏感性\\
Monte Carlo & 采样与方差 & 独立种子、标准误与区间\\
\bottomrule
\end{tabular}
\end{center}

\MFQTransition{例：一层树中的复制}

取 $S_0=100,u=1.2,d=0.8,K=100$，一期无风险总回报为 $R=1.02$。终端看涨支付为 $20,0$。复制要求 $120\Delta+RB=20$、$80\Delta+RB=0$，相减得 $\Delta=1/2$、$B=-40/R$，初始价格为 $50-40/R\approx10.7843$。风险中性概率 $p=(R-d)/(u-d)=0.55$ 给同一个价格 $(0.55\cdot20)/R$。因为 $d<R<u$，两状态概率均严格为正。

CRR 树取 $u=e^{\sigma\sqrt{\Delta t}},d=u^{-1}$ 和
\[
 p=\frac{e^{r\Delta t}-d}{u-d}.
\]
需要先验证 $p\in[0,1]$，再从终端收益反推。256 步价格为 $8.908302$，其误差主要来自时间离散\cite{cox1979option}。

\MFQTransition{从空间导数到三对角系数}

令剩余期限为 $\tau$、$S_j=j\Delta S$，无股息空间算子为 $LV=\tfrac12\sigma^2S^2V_{SS}+rSV_S-rV$。中心差分给
\[
 (L_hV)_j=A_jV_{j-1}+B_jV_j+C_jV_{j+1},
\]
\[
 A_j=\tfrac12\sigma^2j^2-\tfrac12rj,\quad
 B_j=-\sigma^2j^2-r,\quad
 C_j=\tfrac12\sigma^2j^2+\tfrac12rj.
\]
从终端支付 $V_j^0=(S_j-K)_+$ 向增加的 $\tau$ 推进，隐式一步解 $(I-\Delta\tau L_h)V^{n+1}=V^n$；已知边界项移到右端。notebook 会显示一次树的倒推，再比较多层树与闭式价格，并用二分反演隐波。

隐式有限差分在 $S\in[0,S_{\max}]$ 上离散 PDE。每个时间层解三对角线性系统
\[
 (I-\Delta t L_h)V^{n+1}=V^n,
\]
并使用 $V(\tau,0)=0$、$V(\tau,S_{\max})\approx S_{\max}-Ke^{-r\tau}$。事先确定网格给出 $8.908382$。误差同时包含空间离散、时间离散和右边界截断；只加密时间而固定 $S_{\max}$ 不能消除全部误差。

教学实验保持其他维度不变，分别把空间节点减半、时间节点减半，并在保持 $\Delta S$ 近似不变时把右边界从 $400$ 缩到 $300$。相对基准网格的差分别为 $0.020505$、$0.000857$ 与约 $4.4\times10^{-14}$。这些是敏感性诊断，不是严格误差上界；边界差很小只说明本合约与本参数下 $S_{\max}=300$ 已足够远。

\MFQTransition{差分格式与误差来源}

把 PDE 离散为显式、隐式或 Crank--Nicolson 格式。显式法便于理解但受稳定步长约束；隐式法每步解三对角系统，稳定性更好；Crank--Nicolson 二阶时间精度但不光滑 payoff 附近可能振荡，可先做 Rannacher 平滑。

% Generated by tools/build_figures.py; edit figures/figure-specs.json instead.
\MFQFigure{tex/figures/generated/lower-dv-pde.pdf}{有限差分从到期支付出发，沿剩余期限增加的方向求解价格；空间边界和时间步长共同影响数值近似。}{lower-dv-pde}


误差预算至少分为空间截断、时间截断、$S_{\max}$ 边界和线性求解误差。只把网格加密一次不能识别哪一项主导。对无套利价格还应检查非负、随 $S$ 单调、对 $K$ 凸等离散性质。

风险中性 Monte Carlo 使用
\[
 S_T=S_0\exp\left((r-\tfrac12\sigma^2)T+\sigma\sqrt T Z\right)
\]
和贴现收益均值，其中 Brownian 过程记为 $W_t$，且 $Z=W_T/\sqrt T\sim N(0,1)$。20,000 条路径给出 $8.779633$，标准误 $0.096831$；解析值落在约两个标准误内。随机误差按 $N^{-1/2}$ 收敛，方差缩减、准 Monte Carlo 或多层方法改变的是计算效率，不修复模型错设\cite{glasserman2004montecarlo}。

树的步数、PDE 的时间和空间网格、Monte Carlo 的路径数，分别影响不同类型的近似误差。比较价格时，需要同时比较相应的误差大小。解析值偶尔落在模拟置信区间之外，也可能只是区间本身未覆盖真值。

同一计算还可以用以下方法比较：
\begin{itemize}
  \item 闭式用 SciPy \texttt{ndtr}，树用 SciPy 二项分布概率质量；
  \item PDE 用 \texttt{solve\_banded}；Monte Carlo 的独立参照用 SciPy 自适应积分直接积分正态密度下的贴现收益。
\end{itemize}
这些方法使用相同的合约和参数，却采用不同的数值计算方式。确定性方法之间的差异主要来自舍入与离散化；Monte Carlo 与数值积分的差异还包含随机抽样误差，应结合标准误解释。

\section{先反演单点，再校准曲面}

给定单个报价 $C^{\mathrm{mkt}}$，隐含波动率解
\[
 C^{BS}(S,K,r,T,\sigma_{\mathrm{impl}})=C^{\mathrm{mkt}}.
\]
这是一维根求解。每个节点各自找到一个 $\sigma_j$，目标损失接近零是必然结果，不能称为“曲面模型校准”。根必须被静态套利上下界 bracket；Vega 太小时还应报告病态性。

% Generated by tools/build_figures.py; edit figures/figure-specs.json instead.
\MFQFigure{tex/figures/generated/lower-dv-vol-surface.pdf}{隐含波动率同时随执行价和期限变化；切片形状必须与报价单位及无套利条件一起解释。}{lower-dv-vol-surface}


二分法的优点是只依赖单调性和 bracket；Newton 法利用 Vega 收敛更快，但在深度实值、深度虚值或短期限节点可能跳出合法区间。可靠实现可先用套利界拒绝不可能报价，再用带 bracket 的方法求根，同时保存价格残差和 Vega。bid/ask 区间内存在许多可接受隐波时，返回过多小数位只会制造虚假精度。

第二阶段才拟合参数化总方差。事先确定教学模型令 $k=\log(K/F_T)$，
\[
 w(k,T)=\sigma^2(k,T)T=aT+bk^2+cT^2.
\]
对 15 个非等距执行价与期限节点先逐点反演，再以加权最小二乘恢复 $(a,b,c)=(0.035,0.080,0.006)$。这个例子说明报价怎样决定一组共享参数，以及拟合残差怎样反映参数化限制。该形式本身不保证一般的曲面无套利；更深入的研究可参见局部波动率等模型\cite{dupire1994pricing}。

校准试图由多个报价确定一组共享参数。设计矩阵接近奇异时，不同参数可能产生近似相同的样本内价格，却对缺失期限给出不同外推。可以比较参数约束、报价权重和留出节点的重定价误差。若权重来自买卖价差所反映的报价精度，应先定义这一关系；根据拟合残差事后降低某些权重，会改变原先的估计问题。

\MFQTransition{非等距凸性与期限条件}

看涨价格随执行价递减。非等距网格的凸性应比较斜率
\[
 \frac{C_i-C_{i-1}}{K_i-K_{i-1}}
 \le
 \frac{C_{i+1}-C_i}{K_{i+1}-K_i},
\]
不能直接用 $C_{i-1}-2C_i+C_{i+1}\ge0$。后者只在等距网格对应离散二阶差分。

固定执行价的看涨价格随期限不降，需要无分红与非负利率等条件。存在分红、负利率或不同远期时，应先明确贴现因子与远期归一化，再选择适用的日历无套利条件；期限之间的比较因此依赖所采用的贴现、远期与执行价约定。

归一化模式令 $F_T=S_0e^{(r-q)T}$、$D_T=e^{-rT}$，在固定
$k=K/F_T$ 上比较 $C/(D_TF_T)$。教学实验取非零股息率 $q=3\%$，由每个节点携带
\texttt{forward} 与 \texttt{discount\_factor}；网格不完整或原始固定执行价模式与假设冲突时均拒绝。这样“通过日历检查”同时携带远期、贴现和 moneyness 口径。

\MFQTransition{从公开利率到贴现输入}

合成报价节点用于理解参数反演，公开利率例子则演示如何把决策日前可得的利率转换为贴现因子；国债 par yield 也不等于期权定价所需的完整零息曲线。该快照不能验证隐波、校准或交易收益。

从 par yield 到零息贴现因子通常还需 bootstrap，并明确日计数、复利方式和结算日。本例直接使用连续复利，便于理解贴现计算；具体合约仍需要相应的曲线、计息和结算约定。把 1 年 par yield 直接塞进所有期限的 Black--Scholes 公式，会同时制造期限错配与插值偏差。

定价方法近似当前模型下的价格，校准则从报价反推参数。接下来考察这些近似怎样随网格和输入变化，并进一步联系到离散对冲误差。

\MFQTransition{抽样与反演的联系}
欧式看涨的 Monte Carlo 可用 $S_T$ 作控制变量，其风险中性期望为 $S_0e^{(r-q)T}$；路径依赖支付则还要分析路径离散误差。隐波反演把价格误差放大约为 $\Delta\sigma\approx\Delta C/\mathrm{Vega}$，因此 Vega 很小时，价格数值误差与 bid/ask 都可能造成较大的隐波不确定性。

\begin{diagnostic}
校准误差接近零可能只是每个报价都有一个自由参数。模型风险要看参数稳定、插值外推、无套利和对未参与校准报价的表现，而不是训练目标是否为零。
\end{diagnostic}

定价与校准结果会在衍生品路线综合项目中同对冲误差、成本和压力情景合并评价；单点拟合并不等于可实施模型。
